\reviewsection{Media: Access Changes Who Can Enter the Public Record}

In the Module 3 documentary \textit{Wretches and Jabberers}, Tracy Thresher and Larry Bissonnette are autistic adults with limited speech who had been presumed intellectually incapable and excluded from ordinary schooling. After gaining access to typing, they traveled internationally to challenge public assumptions about disability and intelligence \citep{wurzburg2013wretches}. The meaningful change is not that adults gave them competence. An accessible route made their existing authorship harder to dismiss and allowed their messages to change the public conversation. This connects to \citet{norrie2021context}: a device alone does not create agency. The route must remain usable, partners must treat the message as consequential, and support must preserve authorship. In my classroom, I must neither presume incompetence because speech is limited nor treat one method as universal proof. I must provide multiple routes, confirm meaning respectfully, and let student messages revise instruction and support.

\reviewsection{What I Will Carry Into My Classroom}

In my future classroom, I will treat AAC as a continuously revised multimodal system rather than one device. Students need vocabulary for choosing, refusing, commenting, asking, joking, participating academically, expressing sensory and health needs, protecting privacy, and repairing misunderstandings. They need primary and backup options that travel through arrival, instruction, transitions, group work, lunch, dismissal, home, illness, absence, and re-entry. Adults and peers need to model, wait, notice, confirm, and act without requiring performance or taking over the student's voice. The team must also notice when a gesture, facial expression, object, written message, or pause already communicates effectively rather than demanding device repetition to make the message look more conventional \citep{mirenda2015aids}.

My central measure remains the one established in my Module 2 journal, \textit{Who Has More Power After Support?} \citep{garciabautista2026module2journal}. After communication support is introduced, can the student author more messages, enter more relationships, challenge an interpretation, protect a boundary, and influence what happens next? If support reduces visible behavior while adult control remains untouched, communication is still being managed rather than respected. Meaningful communication therefore looks like more than access to words. It looks like a learner being recognized as the author of a message that matters. It looks like partners who can tolerate being corrected and plans that change when the communicator shows that a support is not working. It looks like family knowledge traveling into school and student-authored information traveling home. Most of all, it looks like a student having something important to say, an accessible way to say it, and enough power for the message to change the world around them.
